% !TeX root = ../../main.tex
% =====================================================================
%  Cómo usar este libro
% =====================================================================
\chapter*{Cómo usar este libro}
\addcontentsline{toc}{chapter}{Cómo usar este libro}
\markboth{Cómo usar este libro}{Cómo usar este libro}

\section*{Estructura}

El libro se divide en ocho partes. La primera es \textbf{tronco común}: la lee
todo el mundo, antes de tocar el CAD.
Las partes II a VII son los \textbf{itinerarios de cada subsistema}. La VIII es
\textbf{transversal} --- costes, fabricación, planificación --- y también es de
lectura obligatoria, porque el coste y la fabricabilidad son restricciones de
diseño, no un trámite del final de la temporada.

Las partes no son independientes. Como mínimo:

\begin{itemize}
  \item Dinámica vehicular (parte II) va \emph{antes} que suspensión (parte III):
        define los objetivos que la suspensión tiene que cumplir.
  \item La simulación de vuelta (capítulo \ref{ch:simulacion_vuelta}) es la
        herramienta con la que se justifican las decisiones de aerodinámica
        (parte V) y de relación de transmisión (parte VII).
  \item La refrigeración (parte VI) y la aerodinámica (parte V) comparten el
        mismo aire: los capítulos \ref{ch:lado_aire} y
        \ref{ch:elementos_aerodinamicos} se leen juntos.
\end{itemize}

\section*{Los recuadros}

Todos los capítulos usan los mismos seis recuadros. Debes poder hojear el libro
buscando solo un tipo de recuadro y sacar algo útil.

\begin{objetivos}
  \item Abre cada capítulo. Si al terminar no sabes hacer esto, vuelve atrás.
\end{objetivos}

\begin{clave}
La idea que hay que retener aunque se olvide todo lo demás del apartado.
\end{clave}

\begin{caso}[cómo se usa]
Lo que hicimos nosotros de verdad: la decisión, los datos de partida, el
resultado en pista y qué haríamos distinto. Es la parte del libro que no puedes
encontrar en ningún otro sitio, y por tanto la más importante.
\end{caso}

\begin{regla}[cómo se usa]
Requisito del reglamento que condiciona el diseño, con el artículo citado.
\textbf{Contrasta siempre con la versión vigente}: el reglamento cambia todos
los años y este libro no.
\end{regla}

\begin{ojo}
Los errores que se repiten cada temporada. Se escriben aquí precisamente para
dejar de repetirlos.
\end{ojo}

\begin{entregable}
Lo que produces al terminar el capítulo: plantilla, hoja de cálculo, checklist o
apartado del informe. Sin entregable, el capítulo no se da por leído.
\end{entregable}

\begin{juez}
Preguntas tipo del design event sobre el capítulo. Si no sabes contestarlas, no
has entendido el capítulo, independientemente de lo que creas.
\end{juez}

\section*{Cómo se imparte como curso}

Cada capítulo está pensado como una sesión de 90 minutos con este reparto:

\begin{itemize}
  \item \SI{20}{\minute} de teoría mínima --- solo lo necesario para entender el porqué.
  \item \SI{35}{\minute} sobre nuestro caso real (el recuadro \textsc{caso fusal}).
  \item \SI{30}{\minute} de ejercicio, produciendo el entregable.
  \item \SI{5}{\minute} de \emph{qué no sabemos todavía}: la lista de incógnitas
        abiertas que alimenta los TFG y los proyectos del año siguiente.
\end{itemize}

La sesión la imparten siempre \textbf{dos personas}: quien lo sabe y quien lo
tendrá que impartir el año que viene. Se graba siempre.

\section*{Responsables por capítulo}

\pendiente{Rellena esta tabla. Un capítulo sin responsable no se actualiza nunca.}

\begin{center}
\begin{tabular}{@{}llll@{}}
\toprule
\textbf{Parte} & \textbf{Responsable} & \textbf{Revisor} & \textbf{Últ. revisión} \\
\midrule
I --- Tronco común       & & & \\
II --- Dinámica vehicular & & & \\
III --- Suspensión        & & & \\
IV --- Chasis             & & & \\
V --- Aerodinámica        & & & \\
VI --- Refrigeración      & & & \\
VII --- Transmisión       & & & \\
VIII --- Transversal      & & & \\
\bottomrule
\end{tabular}
\end{center}

\section*{Qué leer además de esto}

Este libro es el registro de nuestras decisiones, no un sustituto de la teoría.
Estos son los libros a los que remite constantemente; conviene que el equipo
tenga al menos un ejemplar de cada uno.

\begin{description}[leftmargin=0pt,labelsep=0pt,style=nextline]
  \item[Dinámica vehicular] \textcite{milliken_rcvd_1995} es la referencia: densa,
    cara y con diferencia el mejor libro del campo. \textcite{gillespie_fundamentals_1992}
    es más asequible para empezar. Para el neumático,
    \textcite{pacejka_tirevehicle_2012}. \textcite{smith_tunetowin_1978} no tiene
    una sola ecuación y sigue explicando la puesta a punto mejor que casi nadie.
  \item[Suspensión y estructuras] \textcite{staniforth_suspension_2006} para la
    parte práctica y \textcite{smith_engineertowin_1984} para materiales y
    uniones. El dimensionado, con \textcite{budynas_shigley_2020}.
  \item[Aerodinámica] \textcite{katz_racecaraero_1995} está escrito justo para
    esto. Los fundamentos, en \textcite{anderson_aerodynamics_2017}; el CFD, en
    \textcite{versteeg_cfd_2007}.
  \item[Refrigeración] \textcite{incropera_heat_2011} para los fundamentos y
    \textcite{hesselgreaves_compact_2017} para intercambiadores compactos, que es
    lo que realmente montamos. \textcite{kays_london_1984} sigue siendo la fuente
    de los datos de superficie.
  \item[Formula Student] \textcite{seward_racecardesign_2014} recorre el coche
    entero con este nivel de detalle y es el mejor punto de partida para un
    rookie. Las notas técnicas de \textcite{optimumg_tips} son cortas y muy
    aplicables.
\end{description}

\section*{Cómo se edita}

\begin{itemize}
  \item Un fichero por capítulo, en \texttt{Capitulos/<parte>/}. No toques los
        ficheros de otros sin avisar.
  \item Mientras escribes un capítulo, descomenta \verb|\includeonly| en
        \texttt{main.tex} y deja solo el tuyo: compila en segundos.
  \item Las imágenes van en \texttt{Imagenes/<parte>/}, el código en
        \texttt{Codigo/}, los planos en PDF en \texttt{Planos/}.
  \item Símbolos: usa siempre las macros de \texttt{Estilo/macros.tex}
        ($\Cl$, $\Kus$, $\NTU$\ldots). Nunca los escribas a mano.
  \item Unidades: siempre con \texttt{siunitx} (\verb|\SI{85}{\celsius}|),
        nunca a pelo.
  \item Lo que falta se marca con \verb|\pendiente{}| o \verb|\falta{}| y sale
        listado al compilar en modo borrador. Nada de dejarlo «para acordarse».
\end{itemize}
